
Let $\tau\in\operatorname{Aut}(\overline{\Q}/\Q)$.
Choose an odd integer $e$ describing its action on all values of
$\eta$ and $\psi$, and put $\nu=\eta^{(e-1)/2}$.
Coefficientwise conjugation gives a form $g^\tau$ with character
$\psi^e$ and good eigenvalues $\eta(p)^e a_p(E)$.
One can reduce this conjugation assertion to integral weight by
multiplying by the ordinary theta series: it has rational
coefficients, and its fourth power has trivial scalar character.
The algebraic divisibility condition at the cusps is precisely the
one in Ramsey's section model, so division after conjugation is
again a half-integral-weight form.
The coefficientwise twist $g\otimes\nu$ has character
$\psi\nu^2=\psi^e$ and exactly these good eigenvalues.
The identities hold because $\psi/\eta$ is quadratic.
The twist is modular at a level still supported on $S_0$; this
follows as well by writing it as the finite additive translation
sum using the Gauss sum of the primitive character inducing $\nu$.

Pass to a common such level. Both forms lie in its full
good-Hecke eigenspace, so the proportionality of its $D$ and $D_0$
functionals gives
\[
 \tau(a(D))
 =\frac{\nu(D)a(D)}{\nu(D_0)a(D_0)}
 =\nu(D/D_0)a(D).                                  
\]
There is no vanishing reference denominator: $D_0$ is prime to
$S_0$ and $a(D_0)=1$. No rationality of the central-value quotient
has been assumed in obtaining this equality.

Write the Dirichlet character as a product of its finite
prime-power factors $\eta_p$. For every $p$ dividing its conductor
choose $u_p\in\Z_p^\times$ with $u_p^2=D/D_0$, and set
\[
                 z_D=\prod_{p\mid\operatorname{cond}(\eta)}
                                      \eta_p(u_p).
\]
Then $z_D$ is a two-power root of unity,
$z_D^2=\eta(D/D_0)$, and
\[
 \tau(z_D)=z_D^e,\qquad
 \nu(D/D_0)=z_D^{e-1}.
\]
Changing the selected odd representative of $e$ has no effect on
the last identity, because $D/D_0$ is a local square at the
conductor. The preceding Galois identity now shows that $a(D)/z_D$ is fixed by
every $\tau$, hence lies in $\Q$.
Different choices of the local square roots change $z_D$ by a
sign, which changes only the sign of this rational number.

Finally, $\psi/\eta=\chi_{-\sigma}$ is trivial on $D/D_0$:
its conductor is supported at two and the ratio is a square there.
Thus $\psi(D/D_0)=z_D^2$. Dividing (E.10) by $z_D^2$ gives
\[
 a(D)=z_Dr_D,\qquad z_D\in\mu_{2^\infty},\qquad r_D\in\Q,
 \qquad r_D^2=\frac{L_*(\sigma D)}{L_*(\sigma D_0)}.       \tag{E.11}
\]
The argument also covers the zero coefficient case. In particular,
every nonzero coefficient in (E.11) has integral normalized
two-adic valuation.

